\section{L4 -- SIMD (Flynn, NEON, MISD)}

\subsection{Flynn's Taxonomy}

\mult{2}

\begin{definition}{The Four Classes}\\
    instruction streams $\times$ data streams:
    \begin{itemize}
        \item \textbf{SISD}: 1 ALU (single core)
        \item \textbf{SIMD}: 1 instr, many lanes
        \item \textbf{MISD}: lockstep / dependability
        \item \textbf{MIMD}: independent cores
    \end{itemize}
\end{definition}

\begin{concept}{SIMD $\to$ SIMT $\to$ SPMD}\\
    Increasingly abstract versions of the same lockstep-on-data idea (hardware $\to$ Nvidia threads $\to$ software pattern).
\end{concept}

\multend

\begin{KR}{Recipe: Classify with Flynn}\\
    Count instruction streams and data streams: 1/1 SISD, 1/many SIMD, many/1 MISD, many/many MIMD.
\end{KR}

\begin{example2}{Ex.\ S6/9 -- SPMD pattern}\\
    What is SPMD and why is it here?
    \tcblower
    One program per UE, unique ID, branch/index by ID. It is the software generalisation of \important{SIMD} (same instruction, different data).
\end{example2}

\subsection{ARM NEON}

\begin{definition}{NEON Registers \& Lanes}\\
    ARM SIMD unit.
    \begin{itemize}
        \item \textbf{Q} (128-bit) = $2{\times}64$ / $4{\times}32$ / $8{\times}16$ / $16{\times}8$
        \item \textbf{D} (64-bit) = half as many lanes
    \end{itemize}
    % INSERT IMAGE: L4 slide 13 -- register split into lanes
\end{definition}

\begin{concept}{NEON Properties}\\
    \begin{itemize}
        \item built around \important{MAC} (multiply-accumulate)
        \item \important{no divide} (reciprocal estimate only)
        \item type name \texttt{<base><bits>x<count>\_t}
    \end{itemize}
\end{concept}

\begin{KR}{Recipe: Decode a NEON Intrinsic}\\
    Pattern \texttt{v <op> q? \_ <type>}:
    \paragraph{op} the operation (\texttt{mul}, \texttt{mla}, \ldots).
    \paragraph{q?} present $\Rightarrow$ 128-bit Q, else 64-bit D.
    \paragraph{type} \texttt{<base><bits>x<count>}: \emph{count} = lanes, \emph{bits} = lane width.
    \paragraph{result} lane-wise op over \emph{count} elements.
\end{KR}

\begin{example2}{Ex.\ S18/19 -- Decode intrinsics}\\
    \texttt{vmulq\_f64}? lanes in \texttt{uint16x4\_t}? \texttt{vmul\_lane\_u16(a,b,L)}?
    \tcblower
    \texttt{vmulq\_f64}: 128-bit lane-wise multiply of \important{2 doubles} $\to$ \texttt{float64x2\_t}. \texttt{uint16x4\_t}: \important{4 lanes}. \texttt{vmul\_lane}: broadcast \texttt{b[L]} to all lanes, multiply by \texttt{a}.
\end{example2}

\begin{example2}{Ex.\ S20 -- Lab: lanes/iterations}\\
    Process $K$ elements with NEON -- how many vector ops?
    \tcblower
    lanes = vector width / element size; iterations $=\lceil K/\text{lanes}\rceil$. E.g.\ \texttt{uint16x8} $\to \lceil K/8\rceil$ vector ops (+ MAC for the filter).
\end{example2}

\subsection{MISD / Lockstep}

\begin{concept}{Dependable Computing}\\
    MISD spends \emph{redundant hardware} on reliability, not speed (opposite goal to SIMD/MIMD).
\end{concept}

\begin{definition}{Lockstep Variants}\\
    \begin{itemize}
        \item tight: 2 CPUs $0$--$2$ clk apart, compare bus $\to$ RESET (e.g.\ ABS)
        \item loose: async, dissimilar impls + acceptance voting
        \item \textbf{2oo3}: one fault tolerated
    \end{itemize}
\end{definition}

\begin{example2}{Ex.\ S21--24 -- Real-world MISD}\\
    Where do you see lockstep in practice?
    \tcblower
    \important{U96 / Zynq ZU3EG}, $2{\times}$Cortex-R5F in tightly-coupled lockstep.
\end{example2}
